\section{Introduction}
\label{sec:intro}

Protein language models such as \esm \citep{lin2023esm2} predict structure and variant fitness without ever seeing a structural label, which strongly suggests their hidden states encode biology. Mechanistic-interpretability work on these models has converged on a recipe borrowed from large language models: train a sparse autoencoder (\sae) on per-residue activations and inventory the resulting features. The flagship \interplm study \citep{simon2024interplm} reports roughly 2{,}500 \sae features per layer with F1 $> 0.5$ against at least one \swissprot concept, against only 46 such units in the dense neuron basis. The headline observation, however, is that \emph{only $\sim$20\% of \sae features map cleanly to any annotated concept}. The remaining \emph{dark features} have repeatedly been described as candidate ``model-generated scientific hypotheses'': features that fire consistently on protein residues but that no biologist has yet named.

\para{Why this matters.} If even a small fraction of dark features encode real biology, they form a new pipeline for discovery, parallel to BLAST-style sequence search and to AlphaFold-driven structural inference. If most of them are noise, the field's enthusiasm for \sae-based ``interpretability'' on \plm{}s is overstated. Either answer is actionable, and to our knowledge the question has not been answered with quantitative tests that go beyond feature--concept correlation.

\para{Gap.} Existing audits of \sae features in \plm{}s rely almost entirely on feature--concept F1. Three weaknesses follow. (i) F1 is a \emph{correlational} test: a feature with high concept F1 may not causally influence the model's predictions, as recent work on antibody \plm{}s has shown for top-$k$ \sae features \citep{haque2026antibody} and as the broader ``interpretability without actionability'' critique argues for LLMs \citep{bricken2026actionability}. (ii) F1 against a finite concept list cannot tell apart noise from concepts the list does not include---in \swissprot's case, sub-domain motifs (\eg the \tgekp linker between consecutive C2H2 zinc fingers, or the \dry and \npxxy motifs internal to class-A GPCRs) are absorbed into much coarser parent features. (iii) Nobody has paired \sae feature inventories with experimental fitness measurements at the per-residue level to test whether dark features track functional sensitivity. Our work fills these three gaps in one study.

\para{Our approach.} We instrument the pretrained \interplm \sae for \esm layer 4 on a 1{,}500-protein random subset of the human \swissprot proteome (422{,}146 residues) and run a three-axis audit of its dark features. We test (i) whether dark features have statistically more structured activation patterns than a random-residue null; (ii) whether they are causally important to \esm's own masked-LM logits via forward-hook ablation paired with a non-activating control; and (iii) whether \gptfive, prompted with 15 sequence windows per feature, can recover a coherent biological hypothesis. We close the loop by extending the concept list and re-checking F1 against the LLM-proposed labels.

\para{Findings, in numbers.} The top 10 dark features under our pipeline produce 10/10 specific, named hypotheses from \gptfive (9 high-confidence). Eight identify the \tgekp inter-finger linker of C2H2 zinc fingers; two identify the \dry and \npxxy motifs of class-A GPCRs. Adding ``Zinc finger'' to the F1 evaluation lifts 7 of the 8 features above F1 $= 0.6$ (max 0.84), confirming the LLM hypotheses against the same evaluation that originally tagged them as dark. \emph{Causal ablation fails for all 30 top dark features}: the median log$_2$-ratio of KL(orig$\Vert$ablated) at activating vs.\ matched control residues is $-3.86$, and zero features survive Benjamini--Hochberg FDR $< 0.10$ \citep{benjamini1995fdr}. Across the broader population of 1{,}207 ``structured dark'' features, even after adding 7 additional concepts, $\sim$97\% remain dark at F1 $\geq 0.5$, consistent with a long tail of sub-domain motifs that no current annotation system encodes.

\para{Contributions.}
\begin{itemize}[leftmargin=*,itemsep=0pt,topsep=0pt]
    \item We provide the first quantitative, three-axis audit of ``dark'' \sae features in a \plm, combining a structured-vs-random null test, causal forward-hook ablation, and LLM-grounded annotation on the same feature set.
    \item We show that the top dark features in \esm encode \emph{sub-domain motifs that \swissprot describes at coarser granularity} (\tgekp, \dry, \npxxy), not unknown biology, and that adding a single concept (``Zinc finger'') re-claims 7 of 8 features at F1 $\geq 0.6$.
    \item We deliver a quantitative negative result for the ``feature--concept F1 implies mechanism'' assumption: none of 30 top dark features pass single-feature ablation under FDR control, and the median KL effect is in fact $2^{-3.86}\times$ \emph{smaller} at activating residues than at matched controls.
    \item We release the audit pipeline (\sae extraction $\rightarrow$ dark filter $\rightarrow$ entropy null $\rightarrow$ ablation $\rightarrow$ \gptfive annotation $\rightarrow$ concept re-check), runnable end-to-end on a single A6000 in $\sim$25 minutes for any \interplm checkpoint, as a reusable tool for future \plm interpretability work.
\end{itemize}

\para{Organization.} \secref{sec:related} positions our work against \plm interpretability and the ``actionability'' critique. \secref{sec:methodology} describes data, model, \sae, and each audit axis. \secref{sec:results} reports the five experiments. \secref{sec:discussion} draws implications for \plm-driven discovery. \secref{sec:conclusion} concludes with concrete next steps.
